\section{Related work, discussion and open problems}\label{sec:related}

\paragraph{Partially observed control and qualitative POMDPs.}
Keeping a stock above a floor under an unknown but fixed parameter is a safety objective in a partially observable system with a finite, static hidden state. Qualitative analysis of POMDPs---almost-sure and sure safety with belief supports as the relevant state---is the setting of Chatterjee, Doyen and Henzinger \cite{CDH2010}; the blind problem on the cube is the sure-safety problem with a trivial observation function, and the feedback game of Section~\ref{sec:observation} is sure safety with the stock as observation. Point-based value iteration \cite{Pineau2003,Shani2013} and its bounds \cite{Lovejoy1991} approximate quantitative POMDP values on finite horizons; Section~\ref{sec:pbvi} places the exact alpha-sets of the earlier editions in this family and gives the horizon after which finite-horizon survival and indefinite survival coincide ($L+1$ at $m=4,6$; $2\lfloor(L-1)/2\rfloor+1$ at $m=5$). Constrained and chance-constrained POMDPs \cite{Altman1999,Poupart2015,Santana2016,Undurti2010,Ono2015} treat the floor as a constraint in expectation or in probability; here the constraint is sure, and the survivable-set antichain is the exact object.

\paragraph{Antichains.}
Antichain algorithms for universality, model checking and synthesis \cite{DeWulf2006,DeWulf2008,DoyenRaskin2009,Filiot2009,Filiot2011} exploit the fact that the relevant sets of states are up- or down-closed and representable by their minimal or maximal elements. The slack game of Section~\ref{sec:thresholds} is of this kind on $\N^S$, where Dickson's lemma \cite{Dickson1913} guarantees finite antichains but descending chains of up-sets need not stabilise; the saturated game plus forward reachability is the two-sided replacement used here.

\paragraph{Extremal combinatorics of the cube.}
Kleitman's diameter theorem \cite{Kleitman1966} bounds the size of a set of given diameter in the hypercube and identifies the extremal configurations (balls for even diameter, doubled balls of a facet for odd); Corollary~\ref{cor:kleitman} and Theorem~\ref{thm:m6} show that the survivable sets of the cube are governed by it, and that the extremal sets are survivable for $m\le6$. Sperner theory \cite{Sperner1927,Engel1997} supplies the language of antichains in the Boolean lattice; the isoperimetric and Dedekind-type results cited in earlier editions are not used.

\paragraph{Duality.}
Theorem~\ref{thm:criterion} is a minimax theorem on a product of a box and a simplex, the finite-cube instance of the measure certificates of \cite{Abaee2026minimax}, which rest on Sion's theorem and Helly's bound on certificate support; the support bound $m+1$ of Section~\ref{sec:criterion} is the same Helly number. Adjustable robust optimisation \cite{WeiZhang2024} asks when adapting to revealed uncertainty improves a linear programme; the feedback thresholds of Section~\ref{sec:observation} are an exact instance of such a gap ($\Fstar<\infty=\Lstar$) in a safety rather than optimisation setting. Certified and exact computation in reachability analysis \cite{Althoff2020,Geretti2020} shares the motivation of Section~\ref{sec:float}.

\paragraph{What has changed and what is open.}
This edition replaces the enumerative case analyses of the earlier versions by one criterion with certificates, from which the four- and five-cube results follow and the six-cube classification is new; the observation model is computed rather than named; the alpha-set evaluations are placed; and the floating-point question is answered on the sets where it matters. Four questions remain.
\begin{openproblem}[larger cubes]
For $m=7$ the pair bound is $\delta_7=4$, the floor family consists of the $128$ radius-two balls ($29$ cells) and these are the Kleitman-extremal sets, so the largest survivable sets are known; for $m=8$ the largest are the $8\cdot128$ doubled radius-two balls ($58$ cells, $\Lstar=1$). What is not known is the rest of the classification: the maximal survivable sets that are not extremal (at $m=6$ there are two further orbits, with the boundary ten-sets the more numerous) and the level census. The maximal cliques of $G_7$ and $G_8$ have not been enumerated here, and the orbit-and-certificate method needs a symmetry-aware clique enumeration to scale.
\end{openproblem}
\begin{openproblem}[uniform horizon]
Every minimal obstruction found for $m\le6$ has $V=2$. Is $V(O)\le2$ for every minimal obstruction and every $m$, so that the horizon $L$ of Corollary~\ref{cor:uniform} is universal? The possible values of $V$ are rationals of bounded denominator (a vertex of a rational polytope with $\pm1$ data), so the question is whether a value in $(2,\tfr52)$ ever occurs.
\end{openproblem}
\begin{openproblem}[feedback at scale]
$\Fstar(\Theta_m)$ is $23$ and $18$ for $m=4,5$. The recursion is exact but exponential in the class size; a closed form, or the asymptotics in $m$, is open, as is the value of information for partial observations of the stock (a floor indicator only).
\end{openproblem}
\begin{openproblem}[formalisation]
Theorem~\ref{thm:criterion} and the game certificates are elementary and finite; their formalisation would make the five- and six-cube classifications verified computations, by reflection on the certificates stored in \texttt{results.json}.
\end{openproblem}

\section*{Acknowledgements}
The author thanks the anonymous referees of the previous edition, whose reports on the census baseline, the observation model, the point-based proposition, the deadline law and the five-cube certification shaped this reconstruction.
